\section{Experimental setup and datasets}
\subsection{Detector modules and acquisition}
Two opposing detector modules comprise silicon photomultipliers (SiPMs) coupled to
lutetium--yttrium oxyorthosilicate (LYSO) crystals of dimensions
$2\times2\times3\,\mathrm{mm}^{3}$. A $^{22}\mathrm{Na}$ source is positioned centrally
between the modules. Each front end provides an energy output after the first amplification
stage and a timing output after a second stage designed for timing measurements.
Four oscilloscope channels record both modules: channels 1 and 2 contain the energy
waveforms, while channels 3 and 4 contain the timing waveforms. The documented acquisitions
use $100\,\mathrm{ns}$ records sampled at $80\,\mathrm{GS\,s}^{-1}$.

The University of Chicago (UC) and Fondazione Bruno Kessler (FBK) front-end boards are
considered separately. Within each board, timing corrections can use either the energy
waveforms or the timing waveforms; energy-based event selection remains common to these
choices. No pooled training, cross-board performance table or direct cross-board plot is
part of the analysis.

\subsection{Acquisitions and scientific roles}
The benchmark datasets identify three independent acquisitions at a SiPM bias of
$48\,\mathrm{V}$ for each board. They are assigned to control, development and blind roles,
as summarized separately in Tables~\ref{tab:uc-acquisitions} and~\ref{tab:fbk-acquisitions}.
The configured true TOF is zero for these centred-source acquisitions.
Control determines the preprocessing and LED selection rules. Development provides all
model training and selection data. The blind acquisition supplies the final test population.
Thus, these roles are not obtained by dividing one acquisition into three random fractions.
\begin{table}[tbp]
\centering
\small
\begin{tabular}{lll}
\toprule
Population & Acquisition & Bias [V] \\
\midrule
Control & 48V-R0 & \multirow{3}{*}{48} \\
Train (development) & 48V-R1 &  \\
Test (blind) & 48V-R2 &  \\
\bottomrule
\end{tabular}
\caption{UC benchmark acquisitions from the 10-03 acquisition directory. Acquisition names denote independent records.}
\label{tab:uc-acquisitions}
\end{table}

\begin{table}[tbp]
\centering
\small
\begin{tabular}{lll}
\toprule
Population & Acquisition & Bias [V] \\
\midrule
Control & 48V-R0 & \multirow{3}{*}{48} \\
Train (development) & 48V-R1 &  \\
Test (blind) & 48V-R2 &  \\
\bottomrule
\end{tabular}
\caption{FBK benchmark acquisitions from the 10-05 acquisition directory. Acquisition names denote independent records; they are not randomized subsets of a common record.}
\label{tab:fbk-acquisitions}
\end{table}


Event counts refer to the final populations available for each waveform mode and window,
after the complete selection described below. The population can change when the mode or
window changes, because an event must satisfy the applicable timing requirements and provide
the requested waveform interval. Within a fixed board, mode and window, models use the same
prepared development and blind populations. Tables~\ref{tab:uc-selected-events} and
\ref{tab:fbk-selected-events} are reserved for the exported counts; no counts from earlier
analyses are inserted.
\begin{table*}[t]
\centering
\small
\begin{tabular}{lllrrr}
\toprule
Mode & Window & Interval [ns] & Control & Train & Test \\
\midrule
\multirow{2}{*}{Timing} & Short & -1 to 2 & --- & --- & --- \\
 & Extended & -2 to 30 & --- & --- & --- \\
\bottomrule
\end{tabular}
\caption{UC selected-event counts by waveform mode and input window. Train denotes the full development population; test denotes blind data. Dashes denote unavailable counts.}
\label{tab:uc-selected-events}
\end{table*}


\begin{table*}[t]
\centering
\small
\begin{tabular}{lll S[table-format=5.0]}
\toprule
Mode & Window [ns] & Population & {Selected events} \\
\midrule
\multirow[t]{6}{*}{Energy}
 & \multirow[t]{3}{*}{Onishi ($-1.5$ to $2$)}
 & Control & 15488 \\
 & & Development & 77252 \\
 & & Blind & 80365 \\
\addlinespace[3pt]
\cdashline{2-4}
\addlinespace[3pt]
 & \multirow[t]{3}{*}{Wide ($-2$ to $30$)}
 & Control & 15482 \\
 & & Development & 77227 \\
 & & Blind & 80346 \\
\addlinespace[3pt]
\hdashline
\addlinespace[3pt]
\multirow[t]{6}{*}{Timing}
 & \multirow[t]{3}{*}{Onishi ($-1.5$ to $2$)}
 & Control & 14107 \\
 & & Development & 70385 \\
 & & Blind & 73227 \\
\addlinespace[3pt]
\cdashline{2-4}
\addlinespace[3pt]
 & \multirow[t]{3}{*}{Wide ($-2$ to $30$)}
 & Control & 14107 \\
 & & Development & 70385 \\
 & & Blind & 73227 \\
\bottomrule
\end{tabular}
\caption{FBK selected event populations by waveform mode and
input window, after event selection, LED coincidence requirements
and window availability.}
\label{tab:fbk-selected-events}
\end{table*}

